\documentclass[11pt]{article}
\usepackage[T1]{fontenc}
\usepackage{lmodern}
\usepackage{amsmath,amssymb}
\usepackage{booktabs}
\usepackage{array}
\usepackage[a4paper,margin=2.0cm]{geometry}
\usepackage[hidelinks]{hyperref}
\newcommand{\doi}[1]{\href{https://doi.org/#1}{\texttt{#1}}}
\setlength{\parskip}{2pt}
\title{\vspace{-1.5cm}\textbf{Model v4 --- field equations for the three papers}\\[2pt]
\large Revised after review: NSGT in Lim form, local couplings, Peltier reciprocity, symbols}
\author{Dr.\ Vipin \\ Gurugram University}
\date{30 October 2026}
\begin{document}
\maketitle
\vspace{-8mm}
\footnotesize
\noindent\textbf{Scope.} The coupled field equations of the three papers, counted, with a source for every
term. The model stops before the boundary conditions, the dispersion determinant and any solution step.
Status marks: [R]~read in the source, including the equation number I can point to; [T]~standard textbook
law; [D]~form fixed by your instruction; [P]~proposed extension with no source, to be derived by the
scholar from the stated law. Equations I could not open are marked as such and the scholar must read them
before quoting a number.

\section*{0. The size-effect convention, one form for all three papers}

Following your ruling, the nonlocal and gradient operators sit where Lim, Zhang and Reddy put them, and
nowhere else: on the inertia and on the elastic strain energy respectively,
\begin{equation}
\bigl(1-l_{\rm nl}^2\nabla^2\bigr)t_{ij}=\sigma_{ij},\qquad
\sigma_{ij}=c_{ijkl}\bigl(1-l_g^2\nabla^2\bigr)\varepsilon_{kl}
-e_{kij}E_k-\beta_{ij}\theta ,
\tag{N1}
\end{equation}
with the equivalent momentum form, obtained by applying the same operator to the balance of momentum,
\begin{equation}
\bigl(1-l_{\rm nl}^2\nabla^2\bigr)\rho\,\ddot u_i=\sigma_{ij,j}.
\tag{N2}
\end{equation}
Equation (N2) is the form we solve with: the operator is on the inertia term, the stress that appears is
the local one, and no derivative of the operator ever multiplies a boundary term. The differential law is
Eq.~(18) of Lim, Zhang and Reddy, with the constitutive boundary conditions Eq.~(19) of the same paper.
I state this carefully because the Elsevier version of that paper was closed to me; the equation numbers
are the ones identified, sentence by sentence, by Barretta and Marotti de Sciarra, \emph{Int.\ J.\ Eng.\
Sci.} \textbf{130} (2018) 187--198, \doi{10.1016/j.ijengsci.2018.05.009}, who write that Lim et al.\
replaced the integral law by ``the differential law Eq.~(18)'' and that the constitutive boundary
conditions Eq.~(19) hold in wave propagation because the fields vanish at infinity. So (N1) and (N2) are
[R] at one remove: the scholar must open \doi{10.1016/j.jmps.2015.02.001} and confirm (18) on the page
before the number enters a manuscript.

\medskip
\noindent\textbf{Modelling choice, applied to all three papers.} The two operators act on the elastic
strain energy only. Every piezoelectric, piezomagnetic, dielectric, magnetic permeability, pyroelectric,
pyromagnetic, thermal-stress and phonon--phason term is local: it depends on the local fields, never on
their Laplacians. This is your correction and it is the right one for a plate. It keeps $D_i$ and $B_i$
second order, so the electric and magnetic problems stay Poisson-like and need no extra conditions at
$z=\pm h$; it keeps the size effect purely mechanical, so the extra conditions are mechanical too; and it
avoids the double counting that appears when a gradient operator is applied to a term that is already a
field gradient.

\medskip
\noindent\textbf{Symbols and units.} One symbol, one meaning, across the three papers.
\begin{center}\footnotesize
\begin{tabular}{p{2.3cm}p{6.4cm}p{4.9cm}}
\toprule
symbol & meaning & unit\\
\midrule
$k=k_r+ik_i$ & wavenumber along $x$, complex; $k_i>0$ is decay & m$^{-1}$\\
$q$ & magnitude of the carrier charge, $1.602\times10^{-19}$; the electron charge is $-q$ & C\\
$h_{kij}$ & piezomagnetic constant (was $q_{kij}$; renamed to free $q$ for charge) & N\,A$^{-1}$m$^{-1}$\\
$b_n$ & electron mobility (was $\mu_n$; $\mu_{ik}$ is kept for magnetic permeability) & m$^2$V$^{-1}$s$^{-1}$\\
$D_n$ & carrier diffusion coefficient & m$^2$s$^{-1}$\\
$v_d$ & drift velocity, $v_d=b_nE_0$ & m\,s$^{-1}$\\
$\beta_{\rm th}$ & thermodiffusion coefficient, $\beta_{\rm th}=-b_nn_0S_n$; negative for $n$-type carriers & m$^{-1}$s$^{-1}$K$^{-1}$\\
$S_n$ & Seebeck coefficient of the electrons & V\,K$^{-1}$\\
$\Pi_{\rm th}$ & Peltier coefficient, $\Pi_{\rm th}=T_0S_n$ & V\\
$\tau_R$ & carrier recombination time (option, $\tau_R\to\infty$ means none) & s\\
$E_g$ & band gap & J\\
$\rho_p$, $\gamma$ & phason mass density, phason friction & kg\,m$^{-3}$, kg\,m$^{-3}$s$^{-1}$\\
$l_{\rm nl}$, $l_g$ & nonlocal and gradient lengths, entering only as $l/h$ & m\\
$\beta_{ij}=c_{ijkl}\alpha_{kl}$, $p_i$, $m_i$ & thermal-stress tensor, pyroelectric vector, pyromagnetic vector & N\,m$^{-2}$K$^{-1}$, C\,m$^{-2}$K$^{-1}$, N\,A$^{-1}$m$^{-1}$K$^{-1}$\\
$\alpha_{ik}$ & magnetoelectric coefficient of the composite & s\,m$^{-1}$\\
$\tau_q$, $\tau_T$ & Moore--Gibson--Thompson (MGT) phase lag, $\tau_T=(2h)^2/(\pi^2\kappa)$ & s\\
\bottomrule
\end{tabular}
\end{center}

\section*{1. Paper A (Sumit) --- piezoelectric semiconductor plate, ZnO type}

\noindent\textbf{Unknowns (5):} $u_x$, $u_z$, $\varphi$, $\Delta n$, $\theta$. \quad
\textbf{Equations (5):} momentum on (N2), Gauss, carrier continuity, MGT heat.
Constitutive side: (N1) with (A2) below and (N2) for momentum.

\noindent\textbf{Constitutive, local couplings.}
\begin{align}
\sigma_{ij}&=c_{ijkl}\bigl(1-l_g^2\nabla^2\bigr)\varepsilon_{kl}-e_{kij}E_k-\beta_{ij}\theta,
\qquad \bigl(1-l_{\rm nl}^2\nabla^2\bigr)t_{ij}=\sigma_{ij}, \tag{A1}\\
D_i&=e_{ikl}\varepsilon_{kl}+\epsilon_{ik}E_k+p_i\theta . \tag{A2}
\end{align}

\noindent\textbf{Field equations.}
\begin{align}
\bigl(1-l_{\rm nl}^2\nabla^2\bigr)\rho\,\ddot u_i&=\sigma_{ij,j}, \tag{A3}\\
D_{i,i}&=-q\,\Delta n, \tag{A4}\\
\frac{\partial\Delta n}{\partial t}&=D_n\nabla^2\Delta n+b_nn_0\nabla\!\cdot\!E
+v_d\frac{\partial\Delta n}{\partial x}+\beta_{\rm th}\nabla^2\theta-\frac{\Delta n}{\tau_R}, \tag{A5}\\
\Bigl(1+\tau_q\frac{\partial}{\partial t}\Bigr)
\Bigl[\rho c_p\dot\theta+T_0\beta_{ij}\dot\varepsilon_{ij}
+T_0p_i\dot E_i+q\Pi_{\rm th}\frac{\partial\Delta n}{\partial t}
+\frac{E_g\Delta n}{\tau_R}\Bigr]&=k_{\rm th}\nabla^2\theta . \tag{A6}
\end{align}

\noindent\textbf{The two reciprocal terms.} The band-gap term $-(E_g/q)\dot{\Delta n}$ is gone. What
replaces it is the Peltier heat of the carrier current, and the pair is fixed by one material function,
the Seebeck coefficient:
\begin{equation}
\beta_{\rm th}=-b_nn_0S_n,\qquad \Pi_{\rm th}=T_0S_n,\qquad
\frac{\Pi_{\rm th}}{\beta_{\rm th}}=-\frac{T_0}{b_nn_0}\quad\text{(Onsager symmetry).}
\tag{A7}
\end{equation}
For $n$-type carriers of charge $-q$ the Seebeck coefficient is negative, $S_n<0$, so
$\beta_{\rm th}<0$: a temperature gradient drives electrons towards the cold face, and $\Pi_{\rm th}<0$
means that heat is released when carriers pile up, which is what $-q\Pi_{\rm th}\partial_t\Delta n>0$
says for $\partial_t\Delta n>0$. The scholar derives (A6) from the entropy balance, writes the entropy
production, and shows that the cross coefficients of (A5) and (A6) are equal, which is the only honest
way to introduce two new coefficients in one paper. The recombination option is a pair: $-\Delta n/\tau_R$
in (A5) and $+E_g\Delta n/\tau_R$ in (A6), both zero as $\tau_R\to\infty$. If the scholar prefers
$\tau_R=\infty$, both terms are deleted and nothing else changes.

\begin{center}\footnotesize
\begin{tabular}{p{2.6cm}p{6.0cm}p{5.2cm}}
\toprule
term & meaning & source (status)\\
\midrule
(A1) NSGT, (A3) & Lim form, operator on inertia and on $c{:}\varepsilon$ only & Lim, Zhang \& Reddy, \emph{J.\ Mech.\ Phys.\ Solids} \textbf{78} (2015) 298, \doi{10.1016/j.jmps.2015.02.001}, Eqs.~(18), (19) \emph{as identified by} Barretta \& Marotti de Sciarra (2018) [R, secondary]\\
$\beta_{ij}\theta$ & thermal stress by Duhamel's analogy & textbook [T]\\
$p_i\theta$ & pyroelectric polarisation & textbook [T]\\
(A4) & Gauss law with the mobile charge & textbook [T]; PSC form: Yang \& Zhou, \emph{Int.\ J.\ Solids Struct.} \textbf{42} (2005) 3171, \doi{10.1016/j.ijsolstr.2004.10.011} [R]\\
(A5) drift, diffusion & published PSC transport terms & Yang \& Zhou (2005) [R]\\
(A5) $\beta_{\rm th}\nabla^2\theta$ & thermodiffusion (Soret) term & textbook irreversible thermodynamics [T]\\
(A6) operator & MGT, base form $K^*=0$, one time derivative cancelled & form fixed by your instruction [D]; attendant paper: Gupta et al., \emph{Int.\ J.\ Numer.\ Methods Heat Fluid Flow} \textbf{35}(9) (2025) 3400, \doi{10.1108/HFF-04-2025-0254}\\
(A6) $-q\Pi_{\rm th}\partial_t\Delta n$ & Peltier heat, Onsager-reciprocal to the Soret term & textbook thermoelectricity [T], Kelvin relation $\Pi=TS$; used in a PSC continuum by Qu, Pan, Zhu, Jin \& Roy, \emph{Int.\ J.\ Eng.\ Sci.} \textbf{182} (2023) 103775, \doi{10.1016/j.ijengsci.2022.103775} [R, abstract and equations as published]\\
(A6) recombination pair & optional, only with finite $\tau_R$ & [P] from the energy balance of the electron gas\\
\bottomrule
\end{tabular}
\end{center}

\noindent\textbf{Dimension test, every term.}
\begin{center}\footnotesize
\begin{tabular}{p{4.6cm}p{3.0cm}p{5.4cm}}
\toprule
term & unit & equals\\
\midrule
$\partial_t\Delta n$, $D_n\nabla^2\Delta n$, $v_d\partial_x\Delta n$, $\Delta n/\tau_R$ & m$^{-3}$s$^{-1}$ & yes\\
$b_nn_0\nabla\!\cdot\!E$: (m$^2$V$^{-1}$s$^{-1}$)(m$^{-3}$)(Vm$^{-1}$m$^{-1}$) & m$^{-3}$s$^{-1}$ & yes\\
$\beta_{\rm th}\nabla^2\theta$: (m$^{-1}$s$^{-1}$K$^{-1}$)(Km$^{-2}$) & m$^{-3}$s$^{-1}$ & yes\\
$\rho c_p\dot\theta$: (kg\,m$^{-3}$)(J\,kg$^{-1}$K$^{-1}$)(Ks$^{-1}$) & W\,m$^{-3}$ & yes\\
$T_0\beta_{ij}\dot\varepsilon_{ij}$: (K)(N\,m$^{-2}$K$^{-1}$)(s$^{-1}$) & W\,m$^{-3}$ & yes\\
$T_0p_i\dot E_i$: (K)(C\,m$^{-2}$K$^{-1}$)(V\,m$^{-1}$s$^{-1}$) & W\,m$^{-3}$ & yes\\
$q\Pi_{\rm th}\partial_t\Delta n$: (C)(V)(m$^{-3}$s$^{-1}$) & W\,m$^{-3}$ & yes\\
$E_g\Delta n/\tau_R$: (J)(m$^{-3}$)(s$^{-1}$) & W\,m$^{-3}$ & yes\\
$k_{\rm th}\nabla^2\theta$: (W\,m$^{-1}$K$^{-1}$)(Km$^{-2}$) & W\,m$^{-3}$ & yes\\
\bottomrule
\end{tabular}
\end{center}

\noindent\textbf{Novelty, stated as a search result.} To our knowledge, after searching Crossref and the
open web in October 2026 for guided waves in piezoelectric semiconductor plates, no paper solves the
dispersion problem of a semiconducting plate in which the carrier field is coupled to a temperature field
that carries a relaxation time. The neighbours and what each lacks: Gupta, Awwad, Ismail, Ahmad and Das, \emph{Int.\ J.\ Numer.\ Methods Heat Fluid Flow} \textbf{35}(9) (2025) 3400, \doi{10.1108/HFF-04-2025-0254}, already put the MGT heat equation into a piezo-semiconductor half space with voids and study the Rayleigh wave, so the heat operator in a semiconductor is theirs and only the finite plate is ours; Othmani, Takali and Njeh,
\emph{Superlattices Microstruct.} \textbf{106} (2017) 86, \doi{10.1016/j.spmi.2017.03.036}, solve PSC plate
Lamb waves with carriers but no temperature field; Qu et al.\ (2023) build the thermoelectric PSC
continuum with Seebeck and Peltier but no wave problem; Sharma, Pal and Chand, \emph{J.\ Therm.\ Stresses}
\textbf{27} (2004) 33, \doi{10.1080/01495730490255709}, solve piezothermoelastic plate Lamb waves with
attenuation and pyroelectricity but no free carriers; Kumar and Kansal, \emph{Int.\ J.\ Solids Struct.}
\textbf{45} (2008) 5890, \doi{10.1016/j.ijsolstr.2008.07.005}, add a diffusing species to a thermoelastic
plate but their third field is not a carrier gas. Our combination is the unoccupied corner of that set,
not a new physics law.

\noindent\textbf{Order of the system.} Because (N1) and (N2) raise the order in $z$, the plate needs more
boundary conditions than the classical six plus electric and thermal ones, and which extra conditions to
use at $z=\pm h$ is your scholar's problem to settle, not part of this model.

\section*{2. Paper B (Neha) --- 1D hexagonal piezoelectric quasicrystal plate}

\noindent\textbf{Unknowns (5):} $u_x$, $u_z$, $w_z$, $\varphi$, $\theta$. \quad
\textbf{Equations (5):} momentum on (N2), phason balance, Gauss, MGT heat.
\begin{align}
\sigma_{ij}&=C_{ijkl}\bigl(1-l_g^2\nabla^2\bigr)\varepsilon_{kl}+R_{ijkl}w_{kl}-e_{kij}E_k-\beta_{ij}\theta,
\qquad \bigl(1-l_{\rm nl}^2\nabla^2\bigr)t_{ij}=\sigma_{ij}, \tag{B1}\\
H_{ij}&=R_{klij}\varepsilon_{kl}+K_{ijkl}w_{kl}-\beta^{w}_{ij}\theta, \tag{B2}\\
D_i&=e_{ikl}\varepsilon_{kl}+\epsilon_{ik}E_k+p_i\theta, \tag{B3}\\
\bigl(1-l_{\rm nl}^2\nabla^2\bigr)\rho\ddot u_i&=\sigma_{ij,j}, \tag{B4}\\
H_{zj,j}&=\rho_p\ddot w_z+\gamma\dot w_z, \tag{B5}\\
D_{i,i}&=0, \tag{B6}\\
\Bigl(1+\tau_q\frac{\partial}{\partial t}\Bigr)
\Bigl[\rho c_p\dot\theta+T_0\beta_{ij}\dot\varepsilon_{ij}
+T_0\beta^{w}_{ij}\dot w_{ij}+T_0p_i\dot E_i\Bigr]&=k_{\rm th}\nabla^2\theta, \tag{B7}
\end{align}
with $w_{ij}=\partial w_i/\partial x_j$. The phason stress (B2) is classical in the main model; the
NSGT-on-phason variant is a labelled sensitivity result, because no published work applies the operator
to phason strain and there is no measured phason correlation length.

\noindent\textbf{The phason equation, and why (B5) is not an invention.} The telegraph form with both
inertia and friction is the standard dynamic phason balance; Agiasofitou and Lazar, \emph{Int.\ J.\
Solids Struct.} \textbf{51} (2014) 923--929, \doi{10.1016/j.ijsolstr.2013.11.016}, write the phason
momentum balance with the friction term $D_{ij}\dot u^{\perp}_j$ in their Eq.~(26) and the resulting
telegraph equation in their Eq.~(30), and they note that the diffusion limit and the wave limit are the
two ends of the same equation. Baggioli and Landry, \emph{SciPost Phys.} \textbf{9} (2020) 062,
\doi{10.21468/SciPostPhys.9.5.062}, write the same crossover as $\omega^2+i\gamma\omega=v^2k^2$ in their
Eq.~(17): diffusive at small $k$, propagating at large $k$. The microscopic origin, in one line, is in
Lubensky, Ramaswamy and Toner, \emph{Phys.\ Rev.\ B} \textbf{32} (1985) 7444,
\doi{10.1103/PhysRevB.32.7444}, whose hydrodynamic treatment is for the \emph{icosahedral} phase; I cite
it for the relaxation mechanism, not for 1D hexagonal constants, and the scholar should not transfer its
moduli.

\begin{center}\footnotesize
\begin{tabular}{p{2.6cm}p{6.0cm}p{5.2cm}}
\toprule
term & meaning & source (status)\\
\midrule
$R_{ijkl}$, $K_{ijkl}$ & phonon--phason coupling and phason moduli, 1D hexagonal & material set of Zhang, Yu, Zhang \& Elmaimouni, \emph{Acta Mech.} \textbf{232} (2021) 135, \doi{10.1007/s00707-020-02811-7} [R, abstract and equations as published]\\
$\beta^{w}_{ij}\theta$ & thermo-phason coupling & [P]: Duhamel's analogy applied to the phason stress; no source found, and the scholar must say so in the paper\\
(B5) inertia $+$ friction & telegraph phason dynamics & Agiasofitou \& Lazar (2014), Eqs.~(26), (30) [R]; Lubensky et al.\ (1985) for the mechanism, icosahedral [R]\\
(B7) $T_0\beta^w_{ij}\dot w_{ij}$ & work of the phason stress against temperature & [P]\\
(B7) operator & MGT base form & [D]\\
\bottomrule
\end{tabular}
\end{center}

\noindent\textbf{Novelty, stated as a search result.} To our knowledge, after searching Crossref and the
open web in October 2026 for Lamb waves in 1D hexagonal piezoelectric quasicrystal plates, no paper
carries both phason terms of (B5) \emph{and} a relaxation-time heat equation in one dispersion problem.
The neighbours are close and must all be read before the claim is written: Zhang et al.\ (2021) for the
elastic guided-wave baseline; Wang, Yu, Zhang, Elmaimouni and Ming, \emph{Wave Motion} \textbf{134} (2025)
103479, \doi{10.1016/j.wavemoti.2024.103479}, for nonlocal (integral) PQC nanoplates; Wang et al.,
\emph{Acta Mech.} \textbf{233} (2022) 3021, \doi{10.1007/s00707-022-03274-8}, for couple-stress PQC
nanoplates; Wang et al., \emph{Compos.\ Struct.} \textbf{370} (2025) 119430,
\doi{10.1016/j.compstruct.2025.119430}, for layered PQC nanoplates. The dangerous one is Zhang, Song, Li,
Han, Yu and Yang, \emph{Acta Mech.\ Solida Sin.} (2026), \doi{10.1007/s10338-026-00755-4}, generalized
thermoelastic Lamb waves in 1D hexagonal QC nanoplates: it already puts heat into a QC plate, and I could
not read its equations. Week one of this paper is to obtain it and write down exactly which heat operator
it uses. If it uses MGT with the same two lags, the novelty of (B7) is gone and what remains is (B5) with
both phason terms.

\noindent\textbf{Order of the system.} (N1) and (N2) raise the order in $z$ on the phonon side as in
Paper A, and (B5) adds a second-order phason equation with no counterpart in a classical plate; how many
extra conditions at $z=\pm h$ the coupled system needs, and which of them are phason conditions, is your
scholar's problem to settle, not part of this model.

\section*{3. Paper C (Anu) --- magneto-electro-elastic plate, BaTiO$_3$/CoFe$_2$O$_4$}

\noindent\textbf{Unknowns (5):} $u_x$, $u_z$, $\varphi$, $\psi$, $\theta$, with $H_i=-\psi_{,i}$. \quad
\textbf{Equations (5):} momentum on (N2), electric Gauss, magnetic Gauss, MGT heat.
\begin{align}
\sigma_{ij}&=C_{ijkl}\bigl(1-l_g^2\nabla^2\bigr)\varepsilon_{kl}-e_{kij}E_k-h_{kij}H_k-\beta_{ij}\theta,
\qquad \bigl(1-l_{\rm nl}^2\nabla^2\bigr)t_{ij}=\sigma_{ij}, \tag{C1}\\
D_i&=e_{ikl}\varepsilon_{kl}+\epsilon_{ik}E_k+\alpha_{ik}H_k+p_i\theta, \tag{C2}\\
B_i&=h_{ikl}\varepsilon_{kl}+\alpha_{ik}E_k+\mu_{ik}H_k+m_i\theta, \tag{C3}\\
\bigl(1-l_{\rm nl}^2\nabla^2\bigr)\rho\ddot u_i&=\sigma_{ij,j}, \tag{C4}\\
D_{i,i}&=0,\qquad B_{i,i}=0, \tag{C5}\\
\Bigl(1+\tau_q\frac{\partial}{\partial t}\Bigr)
\Bigl[\rho c_p\dot\theta+T_0\beta_{ij}\dot\varepsilon_{ij}
+T_0p_i\dot E_i+T_0m_i\dot H_i\Bigr]&=k_{\rm th}\nabla^2\theta . \tag{C6}
\end{align}
The piezomagnetic constant is written $h_{kij}$ so that $q$ stays the carrier charge in Paper A and $k$
stays the wavenumber in all three.

\noindent\textbf{What the neighbours actually do.} Pyroelectric and pyromagnetic coefficients \emph{are}
retained in plate papers, so the earlier draft's claim that they are rarely kept is withdrawn. Kondaiah,
Shankar and Ganesan, \emph{Coupled Syst.\ Mech.} \textbf{2} (2013) 1, \doi{10.12989/csm.2013.2.1.001},
keep both in an MEE plate; Tocci Monaco, Fantuzzi, Fabbrocino and Luciano, \emph{Nanomaterials}
\textbf{11} (2021) 87, \doi{10.3390/nano11010087}, keep both in an MEE nonlocal strain gradient \emph{nano}plate
with hygro-thermal loads; Xiao, Han, Liu and Li, \emph{Compos.\ Struct.} \textbf{153} (2016) 704,
\doi{10.1016/j.compstruct.2016.06.063}, solve guided waves in a functionally graded MEE plate with no
thermal field at all. The earlier claim that $\alpha$ is taken frequency dependent is also withdrawn: I
have no source for it. In the papers I read, $\alpha$ is a constant, either measured for a composite or
computed once by homogenisation. Our treatment, an effective $\alpha$ from two homogenisation schemes
carried as a band, is a modelling choice and is labelled as one.

\noindent\textbf{Novelty, stated as a search result.} To our knowledge, after searching Crossref and the
open web in October 2026 for guided waves in magneto-electro-elastic plates, no paper reports the
dispersion and attenuation of Lamb modes in an MEE plate whose heat conduction carries a relaxation time,
with pyroelectric and pyromagnetic coupling both present and with the mechanical size effect of (N1). The
nearest paper, Xiao et al.\ (2016), is a guided-wave MEE plate with no temperature; the nearest thermal
ones, Kondaiah et al.\ (2013) and Tocci Monaco et al.\ (2021), are static, buckling or free-vibration
problems without the complex wavenumber. The risk is stated plainly: the corner is narrow, and the paper
must show on one figure that the attenuation it predicts is larger than the spread between the two
homogenisation schemes.

\noindent\textbf{Order of the system.} As in the other two papers, (N1) and (N2) raise the order in $z$
and the extra conditions at $z=\pm h$ are your scholar's problem to settle, not part of this model.

\section*{4. Honest note on where the earlier findings still bite}

Two facts from the earlier rounds do not go away and should appear as results. First, in an unbounded
medium both Fourier and MGT conduction make the thermoelastic attenuation saturate; MGT rescales the
plateau by $\tau_T^2/(\tau_T-\tau_q)^2$ and moves the peak from $1/\tau_T$ to $1/|\tau_T-\tau_q|$, and it
does not remove the floor. Second, visibility depends on thickness: the crossing condition is
$(2h)^*=\pi\sqrt{\kappa\tau_q}$, and with literature values for ZnO the crossing sits below the thickness
at which the continuum heat model is defensible. Each paper therefore carries one figure early that maps
thickness against $\tau_q$ and shades the region where relaxation is resolvable. Inside the shaded region
MGT does the work; outside it, the adiabatic and isothermal solutions bracket the answer to a width of
order the thermoelastic coupling, and that bracket is reported as the verification of the full model, not
as its replacement.

\section*{5. Counts, and what is still open}

\begin{center}\footnotesize
\begin{tabular}{lcccl}
\toprule
paper & unknowns & field equations & coupled terms beyond the closest papers & numbers\\
\midrule
A (Sumit) & 5 & 5 & Soret--Peltier pair, pyroelectric, drift, recombination option & (A1)--(A7)\\
B (Neha) & 5 & 5 & telegraph phason, thermo-phason work in the heat equation & (B1)--(B7)\\
C (Anu) & 5 & 5 & pyroelectric with pyromagnetic, effective $\alpha$ as a band, MGT heat & (C1)--(C6)\\
\bottomrule
\end{tabular}
\end{center}

\noindent\textbf{Open, and honest about it.} The Lim, Zhang and Reddy equation numbers come through a
secondary source; the AMSS 2026 thermoelastic QC plate paper could not be read; the NSGT constants
$\tau_q$, $\tau_T$ and the MEE thermal properties rest on literature values collected earlier with the
gaps recorded in the internal logs. None of these blocks the packages, and each is written into the
scholar's first-week tasks instead of being papered over.

\noindent\textbf{Deliberately absent:} boundary conditions, interface and face conditions, the dispersion
determinant, the solution method and every numerical value. Those are the scholars' work.

\end{document}
